\begin{abstract}
Although large unsupervised language models (LMs) acquire extensive factual knowledge and certain reasoning abilities, effectively guiding their outputs remains challenging owing to the fully unsupervised approach taken during pretraining. Current approaches to improving this steerability typically rely on gathering human judgments regarding the relative merits of model outputs and subsequently fine-tuning the base LM to reflect these preferences—commonly through reinforcement learning from human feedback (RLHF). However, RLHF is a multi-step and often brittle method that involves first constructing a reward model aligned with human preferences, followed by reinforcement learning to adapt the initial LM in a way that optimizes this surrogate reward while constraining deviation from the original distribution. In this work, we present a novel formulation for the reward model used in RLHF, enabling analytic recovery of the optimal policy in closed form and allowing the canonical RLHF objective to be addressed with a straightforward classification loss. This method, termed \textit{Direct Preference Optimization} (DPO), yields a robust, efficient, and resource-conserving algorithm that obviates the need for LM sampling during fine-tuning and reduces the requirement for intensive hyperparameter search. Empirical evaluations demonstrate that DPO is capable of tuning LMs to human preferences on par with, or superior to, prevailing techniques. In particular, DPO fine-tuning outperforms PPO-based RLHF in sentiment control tasks and achieves comparable or improved results in summarization and single-turn dialogue benchmarks, all while offering greater ease of implementation and training.
\end{abstract}

\section{Introduction}
Large unsupervised language models (LMs) trained on massive text corpora have demonstrated remarkable and unexpected competencies~\citep{chowdhery2022palm, brown2020language, touvron2023llama,bubeck2023sparks}. Nonetheless, because these models are trained on data authored by humans with diverse intentions, preferences, and abilities, their learned behaviors may sometimes reflect goals or skills that are not always desirable. For instance, while it is beneficial for an AI coding assistant to \textit{recognize} common programming errors to effectively correct them, we would prefer the system’s generated code to emulate the (perhaps infrequent) instances of expert programming that exist in its training set. Analogously, it is valuable for a language model to \textit{know} of a widespread misconception endorsed by half the population, but it would be problematic if the model repeated this falsehood in half of the responses concerning it. Therefore, carefully choosing which \emph{responses and behaviors} are exhibited by the model—out of its extensive \textit{knowledge and abilities}—is essential for ensuring AI systems remain safe, effective, and manageable~\citep{ouyang2022training}. Although current strategies predominantly use reinforcement learning (RL) to adjust LMs in line with human judgments, we demonstrate that the optimization criterion employed in these RL-based methods can be reformulated and solved directly using a simple binary cross-entropy loss, substantially streamlining the preference learning workflow.

Broadly speaking, the typical approach in current systems is to instill desired attributes into a language model by leveraging curated collections of human preferences that highlight behaviors judged as safe and beneficial. This process of learning from preferences is applied after a preliminary phase of large-scale, unsupervised pre-training on broad textual datasets. While one straightforward strategy for preference alignment is to carry out supervised fine-tuning using demonstrations of strong responses produced by humans, the predominant and most effective technique relies on reinforcement learning from human or AI-generated feedback (RLHF/RLAIF; \citep{christiano2017deep,bai2022constitutional}). In RLHF, a reward model is trained using human preference data and subsequently employed to guide the language model policy towards producing outputs assigned high reward, while keeping deviations from the pre-trained model within reasonable limits. Despite the strong results yielded by RLHF—particularly in areas like dialogue and code generation—the methodology introduces significant complexity beyond that of supervised learning, requiring the training of multiple models and iterative sampling within the RL training loop, all of which lead to elevated computational requirements.

This work presents a method for optimizing language models to match human preferences directly, omitting the need for explicit reward modeling or reinforcement learning techniques. We introduce \textit{{Direct Preference Optimization} (DPO)}, an approach that effectively optimizes the same reward-maximization objective with a KL-divergence regularization found in standard RLHF frameworks, but with a design that prioritizes ease of implementation and training. At a high level, DPO adjusts model parameters to increase the log odds of responses favored by humans over those less preferred, while employing an adaptive importance weighting for each example. This adaptive weighting mitigates model collapse, an issue we observe when using naive log probability ratio objectives. Like prior RLHF approaches, DPO is grounded in a formal preference framework—such as the Bradley-Terry model (\cite{bradley1952rankanalysis})—which quantifies the degree of alignment between reward functions and human preference annotations. Distinctively, instead of leveraging this preference model solely to train a reward model before optimizing the language model with respect to it, DPO uses a change-of-variables technique to express the preference loss directly in terms of the policy. As a result, given a collection of human-annotated preferences over model outputs, DPO can train the policy using a straightforward binary cross-entropy loss, yielding an optimal policy relative to an implicit reward shaped by the observed preference data.

Our principal contribution, Direct Preference Optimization (DPO), offers an RL-free alternative for preference-based language model training. Empirical results demonstrate that DPO performs on par with, or surpasses, established approaches including PPO-based RLHF, across tasks such as sentiment control, summarization, and conversational modeling, scaling up to models with 6B parameters.

\section{Related Work}

Increasingly large self-supervised language models have demonstrated the capability to perform certain tasks without direct supervision, utilizing zero-shot \citep{radford2019language} or few-shot prompting methods \citep{gpt3,megatron,chowdhery2022palm}. Nevertheless, their effectiveness on downstream tasks and their alignment with user expectations can be greatly enhanced through fine-tuning on curated datasets containing instructions and human-authored responses \citep{mishra-etal-2022-cross,sanh2022multitask,chung2022scaling,thoppilan2022lamda}. This process, referred to as 'instruction-tuning', equips LLMs with the ability to generalize to instructions that were not explicitly covered during fine-tuning, thereby improving their general applicability and user utility \citep{chung2022scaling}. In practice, acquiring \textit{relative} judgments from humans—where annotators compare responses—is typically more scalable than collecting detailed expert demonstrations. Consequently, recent approaches have shifted toward refining LLMs with preference-based datasets, which has resulted in gains across tasks such as translation \citep{kreutzer-etal-2018-reliability}, summarization \citep{stiennon2022learning,ziegler2020finetuning}, narrative generation \citep{ziegler2020finetuning}, and adherence to instructions \citep{ouyang2022training,ramamurthy2023is}. Typically, these strategies involve first training a reward model—a neural network that estimates response preference—using frameworks like the Bradley-Terry model \citep{bradley1952rankanalysis}, followed by fine-tuning the language model to optimize this reward through reinforcement learning methods such as REINFORCE \citep{williams1992reinforce}, proximal policy optimization (PPO; \cite{schulman2017proximal}), or their extensions \citep{ramamurthy2023is}. A parallel research direction utilizes LLMs already refined with human feedback to create more synthetic preference-labeled data focused on traits like safety or non-harmfulness \citep{bai2022constitutional}, relying primarily on loosely-specified human guidance such as rubrics. These developments lie at the intersection of research on reinforcement learning-based language model training \citep{Ranzato2015SequenceLT,paulus2018a,wu2018learning} and on broader strategies for learning from human preference signals \citep{christiano2017deep,kupcsik2018learning}. However, even with the advantages presented by relative preference data, applying reinforcement learning for the fine-tuning of large language models presents considerable practical barriers; the current work introduces an alternative framework that is theoretically grounded for optimizing relative preference signals, circumventing the need for RL.

Research into learning policies based on preferences, outside of linguistic applications, has been conducted in both bandit and reinforcement learning frameworks, yielding various methodologies. When considering contextual bandit scenarios where preferences or rankings are observed instead of numerical rewards, the problem is formalized as the contextual dueling bandit (CDB) problem \citep{yue2012karmed,dudik2015contextual}. In such frameworks, where explicit rewards are not available, CDB theory replaces the goal of finding an optimal policy with identifying a \textit{von Neumann winner}, defined as a policy that is expected to outperform any alternative policy at least half the time \citep{dudik2015contextual}. Notably, CDB formulations assume that preference feedback arrives sequentially during online learning; by contrast, preference-based policy learning from humans generally utilizes a static, pre-collected dataset consisting of action pairs annotated with offline preferences \citep{yan2022human}. Likewise, in \textit{preference-based RL} (PbRL), algorithms infer policies from binary preference signals determined by an unobserved ‘scoring’ function, not direct rewards \citep{BusaFekete2014,ruiz2023dueling}. PbRL solutions, particularly those leveraging off-policy data, often center on first modeling the underlying scoring (or reward) function and then performing optimization relative to this learned model \citep{jain2013learning,BusaFekete2014,christiano2017deep,sadigh2017active,kupcsik2018learning}. In contrast, we advocate for a unified approach that forgoes explicit reward modeling, focusing instead on optimizing the policy directly to align with observed preferences.

\section{Preliminaries}\label{section:prelims}

Here, we summarize the RLHF pipeline as developed by \citeauthor{ziegler2020finetuning} and later employed in works such as \citep{stiennon2022learning, bai2022training, ouyang2022training}. This process typically involves three main components: (1) supervised fine-tuning (SFT); (2) gathering of preference data and construction of a reward model; and (3) reinforcement learning for policy optimization.

\textbf{SFT}: The initial step in RLHF consists of refining a large, pre-trained language model with supervised training, using high-quality labeled examples relevant to downstream objectives (e.g., conversational agents, text summarization), resulting in the SFT model $\pi^\text{SFT}$.

\textbf{Reward Modelling Phase}: Next, the SFT model is used to respond to prompts $x$, generating answer pairs $(y_1, y_2)\sim \pi^\text{SFT}(y \mid x)$. Human annotators are then presented with these output pairs and asked to indicate a preference, yielding feedback in the form $y_w\succ y_l \mid x$, where $y_w$ is the preferred output and $y_l$ is the less preferred, given $(y_1, y_2)$. This preference data is presumed to arise from some unknown underlying reward function $r^*(y, x)$ that cannot be directly observed. Preferences are commonly modeled probabilistically, with the Bradley-Terry (BT) model \cite{bradley1952rankanalysis} widely used for binary comparisons (with Plackett-Luce models \citep{plackett1975analysis, luce2012individual} applicable when more comprehensive rankings are available). Within the BT framework, the probability of a preference—modeled as the distribution $p^*$—can be formalized as:

\begin{equation}\label{eq:bradley-terry}
    p^*(y_1\succ y_2 \mid x)=\frac{\exp\left(r^*(x, y_1)\right)}{\exp\left(r^*(x, y_1)\right) + \exp\left(r^*(x, y_2)\right)}.
\end{equation}

Given a fixed dataset of pairwise comparisons $\mathcal{D}=\bigl\{x^{(i)}, y_w^{(i)}, y_l^{(i)}\bigr\}_{i=1}^N$ drawn from $p^*$, a reward model $r_{\phi}(x, y)$ can be introduced and its parameters estimated using maximum likelihood techniques. Recasting this setting as a binary classification task, we define the negative log-likelihood loss:

\begin{equation}\label{eq:reward_model}
    \mathcal{L}_R(r_{\phi}, \mathcal{D}) = -\mathbb{E}_{(x, y_w, y_l)\sim \mathcal{D}}\bigl[\log \sigma(r_{\phi}(x, y_w)- r_{\phi}(x, y_l))\bigr]
\end{equation}

with $\sigma$ denoting the logistic function. Within the large language model framework, $r_{\phi}(x, y)$ is commonly initialized from the SFT model $\pi^\text{SFT}(y \mid x)$, and enhanced with a linear projection atop the final transformer block to output a single reward score~\cite{ziegler2020finetuning}. Previous works regularize rewards by enforcing $\mathbb{E}_{x,y\sim \mathcal{D}}\left[r_\phi(x, y)\right] = 0$ across all $x$, thereby reducing variance in the learned reward function.

\textbf{RL Fine-Tuning Phase}: In the reinforcement learning stage, the trained reward function guides model optimization. As in previous studies~\citep{jaques2017sequence, jaques2020human}, the procedure involves solving

\begin{equation}\label{eq:RL}
\max_{\pi_{\theta}}  \mathbb{E}_{x\sim \mathcal{D}, y\sim \pi_{\theta}(y \mid x)}\bigl[r_{\phi}(x, y)\bigr] - \beta\mathbb{D}_{\textrm{KL}}\bigl[\pi_{\theta}(y\mid x)\mid \mid \pi_\text{ref}(y\mid x)\bigr],
\end{equation}

where the hyperparameter $\beta$ modulates the degree to which $\pi_{\theta}$ departs from the reference policy $\pi_\text{ref}$, often corresponding to the SFT model $\pi^\text{SFT}$. The initial state of $\pi_\theta$ is likewise set to $\pi^\text{SFT}$. The KL regularization term constrains the policy to remain close to regions where the reward model remains reliable, and serves to preserve output diversity and avoid collapse to a narrow set of high-reward answers. Since generating language involves discrete decisions, direct optimization is infeasible and reinforcement learning must be employed. Established methods~\citep{ziegler2020finetuning, stiennon2022learning, bai2022training, ouyang2022training} typically define ${r(x, y) = r_{\phi}(x, y) -\beta (\log \pi_{\theta}(y\mid x) - \log \pi_\text{ref}(y\mid x))}$ as the reward, optimizing the policy through PPO~\cite{schulman2017proximal}.

\section{Direct Preference Optimization}\label{sec:DPO}

Given the practical and computational complexities associated with using RL algorithms for large-scale model adaptation, as in language model fine-tuning, we set out to construct a more streamlined methodology for optimizing policies based directly on preference feedback. Departing from the standard RLHF pipeline, which first fits an explicit reward and then maximizes it via RL, our approach harnesses a specific parameterization of the reward function that allows its optimal policy to be derived in closed form—obviating the need for reinforcement learning steps. We now elaborate on this central insight: by exploiting an analytic mapping that links reward models to their associated optimal policies, we can convert a loss on reward functions into a loss on policies themselves. This change-of-variables perspective negates the requirement for a distinct, fitted reward model, while still adhering to well-established human preference models like the Bradley-Terry framework. Essentially, the policy network is tasked with embodying both the generative language

\textbf{Reformulating the DPO objective.} We begin from the same reinforcement learning objective introduced in prior literature, Eq.~\ref{eq:RL}, considering a general reward function $r$. Building upon established results~\citep{peters2007reinforcement, peng2019advantage, korbak2022reinforcement, go2023aligning}, it follows directly that the solution maximizing reward under a KL-divergence constraint, as formulated in Eq.~\ref{eq:RL}, yields:

\begin{equation}\label{eq:op_policy}
    \pi_r(y\mid x) = \frac{1}{Z(x)}\pi_\text{ref}(y\mid x)\exp\left(\frac{1}{\beta}r(x, y)\right),
\end{equation}

where the normalization constant $Z(x) =\sum_{y}\pi_\text{ref}(y\mid x)\exp\left(\frac{1}{\beta}r(x, y)\right)$, known as the partition function, ensures a valid probability distribution. A thorough derivation can be found in Appendix \ref{app:derivation1}. Notably, direct computation of $Z(x)$, even when the reward model $r_{\phi}$ approximates the true reward $r^*$ via MLE, remains computationally intensive \citep{korbak2022reinforcement, go2023aligning}, thus hindering practical implementation of this form. Nonetheless, it is possible to rearrange Eq.~\ref{eq:op_policy} to rewrite the reward in terms of the optimal policy $\pi_r$, the reference policy $\pi_\text{ref}$, and $Z(\cdot)$, whose value remains unknown. By applying a logarithm to both sides of Eq.~\ref{eq:op_policy} and algebraically manipulating the resulting expression, we obtain:

\begin{equation}\label{eq:main_eq}
    r(x,y) =\beta \log \frac{\pi_r(y\mid x)}{\pi_\text{ref}(y\mid x)} + \beta \log Z(x).
\end{equation}

This parameterization can similarly be applied to the true reward function $r^*$ and the associated optimal policy $\pi^*$. Importantly, the Bradley-Terry framework only depends on reward differences between outputs, i.e., ${p^*(y_1 \succ y_2 \mid x) = \sigma(r^*(x, y_1) - r^*(x, y_2))}$. When we substitute the reward parameterization from Eq.~\ref{eq:main_eq} for $r^*(x,y)$ into the preference formulation in Eq.~\ref{eq:bradley-terry}, the normalization factors $Z(x)$ cancel, which allows human preference probabilities to be characterized solely in terms of $\pi^*$ and $\pi_\text{ref}$. Thus, under the Bradley-Terry model, the optimal RLHF policy $\pi^*$ fulfills the preference-based relation:

\begin{equation}\label{eq:objective}
    p^*(y_1\succ y_2 \mid x)=\frac{1}{1 + \exp\left(\beta \log \frac{\pi^*(y_2\mid x)}{\pi_\text{ref}(y_2\mid x)} - \beta \log \frac{\pi^*(y_1\mid x)}{\pi_\text{ref}(y_1\mid x)}\right)}
\end{equation}

A full derivation appears in Appendix~\ref{app:derivation2}. While Eq.~\ref{eq:objective} specifically adopts the Bradley-Terry model, analogous derivations can be constructed for broader families, including the Plackett-Luce models~\citep{plackett1975analysis, luce2012individual}, as detailed in Appendix~\ref{app:plackett_luce_models}.

With an expression for human preference probabilities now dependent only on the optimal policy, we may directly define a maximum likelihood objective for a parameterized policy $\pi_\theta$. Following a similar approach to reward-based modeling (cf. Eq.~\ref{eq:reward_model}), the policy learning objective becomes:

\begin{equation}\label{eq:optimum_model}
    \mathcal{L}_\text{DPO}(\pi_{\theta}; \pi_\text{ref}) = -\mathbb{E}_{(x, y_w, y_l)\sim \mathcal{D}}\left[\log \sigma \left(\beta \log \frac{\pi_{\theta}(y_w\mid x)}{\pi_\text{ref}(y_w

In this formulation, an implicit reward function is modeled via an alternative parametrization such that the optimal policy is identified as $\pi_\theta$. Additionally, this approach is analogous to training a reparametrized Bradley-Terry model, which brings with it desirable theoretical attributes—including consistency—provided certain assumptions about the distribution of preference data hold \cite{bong2022generalized}. We elaborate further on the theoretical aspects of DPO and its connections to related literature in Section~\ref{sec:theory}.

\textbf{How does the DPO update operate?} To better understand the mechanism underlying DPO, one can study the gradient of the loss function $\mathcal{L}_\text{DPO}$. The partial derivative with respect to the parameter vector $\theta$ is given by:

\begin{multline*}\label{eq:gradient}
    \nabla_\theta \mathcal{L}_\text{DPO}(\pi_\theta;\pi_\text{ref}) = \\ -\beta\mathbb{E}_{(x, y_w, y_l) \sim \mathcal{D}} \left[\underbrace{\sigma(\hat{r}_\theta(x, y_l) - \hat{r}_\theta(x, y_w))}_\text{greater impact when reward estimates are inaccurate}\left[\underbrace{\nabla_\theta\log \pi(y_w \mid x)}_\text{encourage $y_w$} - \underbrace{\nabla_\theta\log\pi(y_l \mid x)}_\text{discourage $y_l$}\right]\right],
\end{multline*}

where the implicitly defined reward $\hat{r}_\theta(x, y) = \beta \log \frac{\pi_\theta(y \mid x)}{\pi_\text{ref}(y \mid x)}$ is derived from the language model $\pi_\theta$ and a reference model $\pi_\text{ref}$ (see further discussion in Section~\ref{sec:theory}). At an intuitive level, the loss function gradient acts to increase the probability assigned to more preferred responses $y_w$ while suppressing that of less preferred ones $y_l$. Crucially, the influence of each training example is modulated by the extent to which the implicit reward model $\hat{r}_\theta$ mistakenly assigns higher reward to the less favored completion, multiplied by the factor $\beta$—thus reflecting the degree to which the current model misranks outputs, while respecting the strength of the KL regularization. Our empirical findings indicate that this specific weighting is critical: omitting the weighting term (i.e., employing a more naïve approach) can result in degenerate behavior of the language model (as shown in Appendix Table~\ref{tab:unlikelihood_generations}).

\textbf{DPO overview.} The standard DPO workflow proceeds as follows: First, for each prompt $x$, generate completions $y_1, y_2$ using $\pi_\text{ref}(\cdot \mid x)$, then assign human preference labels to form the offline preference dataset $\mathcal{D} = \{x^{(i)}, y_w^{(i)}, y_l^{(i)}\}_{i=1}^N$. Second, update the language model $\pi_\theta$ by minimizing $\mathcal{L}_\text{DPO}$ given $\pi_\text{ref}$, the dataset $\mathcal{D}$, and a specified $\beta$. In real-world scenarios, practitioners often wish to leverage existing preference datasets instead of generating new samples or conducting additional human preference collection. Because such datasets are generally created using $\pi^\text{SFT}$, the typical practice is to set $\pi_\text{ref} = \pi^\text{SFT}$ when possible. If $\pi^\text{SFT}$ is unavailable, we initialize $\pi_\text{ref}$ by maximizing the likelihood over the preferred completions ${(x, y_w)}$, formally, ${\pi_\text{ref} = \argmax_{\pi}\mathbb{E}_{x, y_w \sim \mathcal{D}}\left[\log \pi(y_w \mid x)\right]}$. This initialization addresses the distribution mismatch between the inaccessible true reference distribution and the $\pi_\text{ref}$ adopted in DPO. Comprehensive implementation details and hyperparameter choices are provided in Appendix~\ref{app:implementation}.

\section{Theoretical Analysis of DPO} In this section, we further elucidate the DPO framework, provide theoretical justification, and discuss how DPO overcomes certain limitations observed in actor-critic approaches commonly applied to RLHF (e.g., PPO~\cite{schulman2017proximal}).

\label{sec:theory}

\subsection{Your Language Model Is Secretly a Reward Model} DPO unifies the processes of reward model estimation and policy optimization, achieving both via a single maximum likelihood objective rather than explicit reward fitting or reinforcement learning steps. The core optimization in Eq. \ref{eq:main_eq} mirrors a Bradley-Terry model, with rewards expressed as $r^*(x, y) = \beta \log\frac{\pi^*_\theta(y \mid x)}{\pi_\text{ref}(y \mid x)}$. The objective is then to optimize the parameterized model $\pi_{\theta}$, which, under a suitable variable transformation, is equivalent to optimizing the reward model as outlined in Eq. \ref{eq:reward_model}. In what follows, we develop the theoretical foundation underlying this reparameterization, demonstrate that it does not limit the expressiveness of the reward models learned, and show that it enables exact identification of the optimal policy. We start by introducing an equivalence relation among reward functions.

\begin{definition}
Two reward functions $r(x, y)$ and $r'(x, y)$ are said to be equivalent if and only if ${r(x, y)-r'(x, y) = f(x)}$ holds for some function $f$.
\end{definition}

It can be readily verified that this definition constitutes an equivalence relation, which results in a partitioning of the set of reward functions into equivalence classes. Based on this, we can state two lemmas:

\begin{lemma}\label{lemma:same_prefrence} Within the framework of Plackett-Luce—particularly, Bradley-Terry—models of preferences, any pair of reward functions belonging to the same equivalence class will produce identical preference distributions.
\end{lemma}

\begin{lemma}\label{lemma:same_policy}
    For any two reward functions from a shared equivalence class, the induced optimal policy in the constrained RL setting is the same.
\end{lemma}

Since their proofs are direct, we postpone them to Appendix \ref{app:lemma1}. The first lemma reflects a well-established identifiability issue in the Plackett-Luce model family \cite{plackett1975analysis}. Owing to this ambiguity, it becomes necessary to enforce extra identifiability constraints in order to provide assurances about the maximum likelihood estimates defined in Eq. \ref{eq:reward_model} \cite{bong2022generalized}. The second lemma establishes that the entire set of reward functions within a given equivalence class results in the same optimal policy, implying that in practice, it suffices to recover any representative from the optimal equivalence class. In Appendix~\ref{app:thm1}, we establish the following theorem:

\begin{theorem}\label{thm:main}
    Assuming certain regularity conditions, every reward class compatible with the Plackett-Luce (and specifically, the Bradley-Terry) model admits the form ${r(x, y) = \beta \log \frac{\pi(y\mid x)}{\pi_\text{ref}(y\mid x)}}$ for some model $\pi(y\mid x)$ and a specified reference model $\pi_\text{ref}(y \mid x)$.
\end{theorem}

\begin{sproof}
    Given any reward function $r(x, y)$, let $\pi_r(y \mid x)$ denote the optimal model it induces as defined by Eq. \ref{eq:op_policy}. We show that there exists a reward function in the equivalence class of $r$ which takes the desired reparameterized form. Define the following projection:
\begin{equation}
    f(r; \pi_\text{ref}, \beta)(x, y) = r(x, y) - \beta\log\sum_{y}\pi_\text{ref}(y\mid x)\exp\left(\frac{1}{\beta}r(x, y)\right)
\end{equation}
This operator $f$ effectively adjusts the reward by subtracting the log partition function over $\pi_r$ for each $x$, and because the term being subtracted depends only on $x$, $f(r; \pi_\text{ref}, \beta)(x, y)$ remains in the equivalence class of $r(x, y)$. Substituting $r$ on the right-hand side of Eq.~\ref{eq:main_eq} (which is valid for any reward function) gives $f(r; \pi_\text{ref}, \beta)(x, y) = \beta \log \frac{\pi_r(y\mid x)}{\pi_\text{ref}(y\mid x)}$. Therefore, $f$ constructs an element in $r$'s equivalence class with the requisite structure, implying that this reparameterization introduces no loss of generality to our reward representation.
\end{sproof}

An alternative interpretation of Theorem~\ref{thm:main} is that it precisely identifies, within each equivalence class of reward functions, the specific form chosen by the DPO reparameterization. Specifically, it chooses the reward function for which

\begin{equation}\label{eq:lag_p}
     \sum_{y}\underbrace{\pi_\text{ref}(y\mid x)\exp\left(\frac{1}{\beta}r(x, y)\right)}_{=\pi(y\mid x)\text{, using Thm.~\ref{thm:main} reparam.}} = 1,
\end{equation}

which guarantees that $\pi(y\mid x)$ forms a legitimate probability distribution, meaning its elements are non-negative and sum to one. Looking at Eq.~\ref{eq:op_policy}, it becomes clear that Eq.~\ref{eq:lag_p} defines the partition function corresponding to the optimal policy derived from the reward function $r(x, y)$. The central idea of the DPO algorithm is to enforce constraints on the typically under-determined Plackett-Luce (and, by extension, Bradley-Terry) preference model family. This approach retains the original expressiveness of the reward models, while also ensuring that the optimal policy described in Eq.~\ref{eq:op_policy} becomes analytically manageable for any input $x$.

\subsection{Instability of Actor-Critic Algorithms} The framework presented here can also illuminate why commonly used actor-critic approaches for RLHF, like PPO, may exhibit unstable behaviors. Focusing on the RL fine-tuning process described in Section \ref{section:prelims}, we can link this setting to the control-as-inference perspective \cite{levine2018reinforcement} in the context of the constrained RL task outlined in \ref{eq:RL}. Suppose we parameterize the policy as $\pi_{\theta}(y\mid x)$ and aim to minimize the KL-divergence $\mathbb{D}_{\text{KL}}[\pi_{\theta}(y|x) \mid \mid \pi^*(y\mid x)]$, where $\pi^*$, defined in Eq.~\ref{eq:optimum_model}, is the optimal policy derived from the reward $r_{\phi}(y, x)$. Through elementary algebra, the optimization objective is given by:

\begin{equation}\label{eq:AC}
    \max_{\pi_{\theta}}\mathbb{E}_{\pi_{\theta}(y\mid x)}\bigg[\underbrace{r_{\phi}(x, y) -\beta\log\sum_{y}\pi_\text{ref}(y\mid x)\exp\left(\frac{1}{\beta}r_{\phi}(x, y)\right)}_{f(r_{\phi}, \pi_\text{ref}, \beta)} - \underbrace{\beta\log\frac{\pi_{\theta}(y\mid x)}{\pi_\text{ref}(y\mid x)}}_{\text{KL}}\bigg]
\end{equation}

This objective aligns with those optimized in earlier studies 
\citep{ziegler2020finetuning, stiennon2022learning, bai2022training, ouyang2022training} by employing the DPO-equivalent reward from the reward class $r_{\phi}$. Under this framework, the normalization term found in $f(r_{\phi}, \pi_\text{ref}, \beta)$ may be viewed as representing the soft value function associated with the reference policy $\pi_\text{ref}$. Although this component does not influence the ultimate optimal solution, its absence can lead to high variance in the policy gradient estimates, thereby undermining training stability. One solution is to approximate the normalization term using a value function trained alongside the policy; however, this approach can pose significant optimization challenges. Another strategy used in previous work involves normalizing rewards via a human completion baseline, which effectively acts as a single-sample Monte Carlo approximation of the normalization term. Distinctively, the DPO reparameterization introduces a reward function that obviates the necessity for such baselines.

\section{Experiments}
Here, we empirically investigate the effectiveness of DPO for policy optimization driven solely by preference data. We begin with a controlled text-generation experiment to examine how effectively DPO balances reward maximization against KL-divergence minimization relative to established preference learning algorithms like PPO. Subsequently, we explore DPO’s utility in scaling to larger model architectures and tackling more challenging RLHF scenarios, such as dialogue and summarization. Our experiments reveal that, with minimal hyperparameter adjustment, DPO matches or surpasses robust baselines including RLHF with PPO and the selection of the best out of $N$ sampled trajectories based on a learned reward model. We first outline our experimental methodology below; further experimental details are provided in Appendix~\ref{app:exp_details}.

\textbf{Tasks.} Our study examines three distinct tasks related to open-ended text generation. Across all experimental conditions, learning agents optimize a policy using a dataset of preference annotations denoted by $\mathcal{D}=\bigl\{x^{(i)}, y_w^{(i)}, y_l^{(i)}\bigr\}_{i=1}^N$. In the \textbf{controlled sentiment generation} task, $x$ represents an initial fragment from a movie review within the IMDb dataset \cite{maas-EtAl:2011:ACL-HLT2011}, and the policy’s objective is to produce continuations $y$ that are positively valenced. To facilitate objective measurement, preference data for this task is synthetically generated by applying a sentiment classifier, ensuring that $p(\text{positive}\mid x,y_w)>p(\text{positive}\mid x,y_l)$. For supervised fine-tuning (SFT), we adapt GPT-2-large until it converges using the training partition of the IMDb dataset (additional experimental settings are described in App~\ref{app:sentiment_details}). For the \textbf{summarization} task, $x$ consists of Reddit forum posts, and the policy generates summaries $y$ that distill the primary content of the post. In line with earlier research, we make use of the Reddit TL;DR summarization corpus \citep{volske-etal-2017-tl} complemented by human preference labels curated by \citeauthor{stiennon2022learning}. The SFT baseline is trained on summaries authored by humans\footnote{\url{https://huggingface.co/CarperAI/openai_summarize_tldr_sft}}, utilizing the TRLX \citep{leandro_von_werra_2023_7790115} infrastructure for RLHF. The human feedback dataset stems from \citeauthor{stiennon2022learning}, who collected judgments on generations produced by a related but separately-trained SFT model. Lastly, in the \textbf{single-turn dialogue} scenario, $x$ is any user-initiated prompt—ranging from factual questions in astrophysics to personal advice requests. The corresponding policy’s role is to deliver a response $y$ that is both informative and engaging; we draw on the Anthropic Helpful and Harmless dialogue dataset \citep{bai2022training}, which includes 170,000 dialogues between users and an AI assistant. Each dialogue concludes with two candidate responses from a large, unspecified language model and an accompanying label that reflects human preference. In this task, a dedicated SFT model does not exist a priori; thus, we construct the SFT baseline by fine-tuning an existing language model on the subset of responses marked as preferred.

\textbf{Evaluation.} Our study adopts two primary evaluation strategies. To rigorously examine each algorithm’s capacity to optimize the constrained reward maximization objective, we focus initially on controlled sentiment generation, where we measure performance along a tradeoff curve between the obtained reward and KL-divergence relative to the reference policy. This frontier is computable in this controlled setting because the exact reward function (a sentiment classifier) is accessible. Conversely, in practical applications where the reward function is unknown, we assess models based on their \textit{win rate} against a designated baseline, utilizing GPT-4 as a surrogate for human judges to evaluate summary quality and response helpfulness in summarization and single-turn dialogue, respectively. In the summarization task, the baseline consists of reference summaries from the test set, while for dialogue, the baseline is the preferred human response from the test examples. While previous work indicates language models may serve as more effective evaluators than established metrics \citep{Chen2023ExploringTU}, we validate our reliance on GPT-4 by conducting a human evaluation, as detailed in Sec.~\ref{sec:human-judgments}. Our findings indicate that GPT-4’s assessments align closely with human judgments, and the rate of agreement between human raters and GPT-4 is on par with, or exceeds, the agreement observed among human annotators.

\textbf{Methods.} In addition to DPO, we benchmark several approaches previously proposed for training language models to match human preferences. For a baseline, we use zero-shot prompting with \textbf{GPT-J} \citep{gpt-j} in the summarization task and a 2-shot prompt for \textbf{Pythia-2.8B} \citep{biderman2023pythia} in the dialogue task. We also evaluate a \textbf{SFT} model and a variant called \textbf{Preferred-FT}, which is produced by fine-tuning on the favored completions $y_w$—drawn either from SFT outputs (for controlled sentiment and summarization) or from a generic language model (for dialogue)—using supervised learning. An additional supervised-inspired approach is \textbf{Unlikelihood}~\citep{welleck2019neural}, which trains the model to increase likelihood on $y_w$ while simultaneously decreasing likelihood on $y_l$, optionally modulated by a coefficient $\alpha\in[0,1]$ controlling the strength of the ‘unlikelihood’ penalty. We further include \textbf{PPO} \citep{schulman2017proximal}, using a reward model trained on preference annotations, and \textbf{PPO-GT}, an oracle variant that employs the exact reward function in the controlled sentiment context. For sentiment experiments, two PPO-GT configurations are used: a standard implementation \cite{leandro_von_werra_2023_7790115}, and a modified variant that applies reward normalization and additional hyperparameter tuning to boost performance (these enhancements are also applied to standard PPO with learned rewards). Finally, we evaluate the \textbf{Best of $N$} strategy, in which $N$ responses are sampled from the SFT model (or Preferred-FT for dialogue), and the completion with the highest reward—according to a model trained from preferences—is selected. Although this method can achieve strong results by divorcing reward model quality from PPO training, it is computationally demanding at inference, as it necessitates generating $N$ completions per prompt.

\subsection{How well can DPO optimize the RLHF objective?}

In RLHF frameworks, the KL-constrained reward maximization objective is designed to strike a balance between maximizing reward and ensuring the learned policy remains close to the reference policy. Consequently, any comparison among algorithms should jointly consider both the attained reward and the KL divergence; achieving marginally greater reward at the expense of a substantially larger KL is often undesirable. Figure~\ref{fig:frontier-tldr-main} illustrates the tradeoff between reward and KL divergence—the so-called frontier—for a range of algorithms evaluated under the sentiment configuration. For each method, we conduct several independent training runs, each employing a distinct value for the policy conservativeness hyperparameter (target KL $\in\{3,6,9,12\}$ for PPO, $\beta \in \{0.05,0.1,1,5\}$, and $\alpha\in\{0.05,0.1,0.5,1\}$ for unlikelihood, and different random seeds for preferred-FT). In total, this produces 22 experimental runs. At intervals of every 100 training steps until the algorithms converge, we assess each resulting policy using a set of held-out prompts, calculating the mean reward given the true reward function, as well as the mean sequence-level KL divergence\footnote{Specifically, this refers to the aggregate of KL-divergences at each individual timestep.} with respect to the reference policy, $\text{KL}\left(\pi\mid \mid \pi_\text{ref}\right)$. Our findings indicate that DPO generates a markedly superior frontier, achieving the greatest reward while maintaining a low KL divergence. This outcome stands out for several reasons. Firstly, although both DPO and PPO optimize an identical objective, DPO exhibits significantly higher efficiency—the reward/KL profile of DPO universally outperforms that of PPO. Secondly, DPO surpasses PPO even in settings where PPO has access to ground truth rewards (PPO-GT), yielding a superior frontier in those cases as well.

\subsection{Can DPO scale to real preference datasets?}
\label{sec:dpo-real-datasets}
We further assess the fine-tuning effectiveness of DPO on tasks involving summarization and single-turn dialogue. In the context of summarization, it is well-established that commonly used automated metrics such as ROUGE exhibit only weak alignment with human judgment~\citep{stiennon2022learning}, motivating prior studies to employ PPO-based fine-tuning using human-labeled preferences to produce higher quality summaries. To compare different fine-tuning approaches, we generate completions on the test set of the TL;DR summarization dataset and determine their average win rate when compared to test reference completions. For all evaluated methods, completions are sampled at a range of temperatures between 0.0 and 1.0; the resulting win rates are illustrated in Figure~\ref{fig:frontier-tldr-main} (right). DPO, PPO, and Preferred-FT were all applied to fine-tune an identical GPT-J SFT model\footnote{\url{https://huggingface.co/CarperAI/openai_summarize_tldr_sft}}. Our experiments indicate that DPO attains a win rate of roughly 61\% when sampling at temperature 0.0, outperforming PPO, which reaches a win rate of approximately 57\% at the same optimal temperature. Additionally, DPO outperforms the best-of-$N$ baseline with respect to peak win rate. It is important to highlight that the $\beta$ hyperparameter for DPO was not systematically optimized; thus, reported results may not represent the upper bound of DPO's capabilities. DPO is observed to exhibit far greater resilience to changes in sampling temperature, in contrast to PPO, whose performance at higher temperatures deteriorates towards that of the original GPT-J model. Notably, Preferred-FT does not provide substantial gains over the SFT baseline. A direct human evaluation comparing DPO and PPO is presented in Section~\ref{sec:human-judgments}; here, DPO outputs generated at a temperature of 0.25 were favored 58\% of the time over PPO outputs sampled at temperature 0.

For single-turn dialogue evaluation, we assess the performance of various approaches on a subset of the Anthropic HH test set \citep{bai2022training} comprising single human-assistant exchanges. The GPT-4-based assessment leverages the annotated preferred completions in the test data as a ground truth to calculate win rates for each method under consideration. Since no standard SFT model exists for this specific task, our procedure begins with the pre-trained Pythia-2.8B model; we subsequently apply Preferred-FT, using the selected completions as training targets to align the reference model's output distribution with the data, followed by DPO training. For comparison, we report results against the strongest out of 128 Preferred-FT generated completions (since the Best of $N$ baseline saturates at 128, as demonstrated in Appendix Figure~\ref{fig:best-of-n}), and also evaluate a two-shot prompted Pythia-2.8B base model. Our findings indicate that DPO performs on par with or exceeds other methods across optimal temperature settings. Additionally, we analyze an RLHF model utilizing PPO and trained on the Anthropic HH dataset\footnote{\url{https://huggingface.co/reciprocate/ppo_hh_pythia-6B}}, sourced from a reputable implementation\footnote{\url{https://github.com/CarperAI/trlx/tree/main/examples/hh}}, but are unable to identify prompts or sampling temperatures that yield results superior to the unmodified Pythia-2.8B. Given our TL;DR results and the shared reward optimization objective of both approaches, we treat the Best of 128 baseline as an approximate indicator of PPO performance. Notably, DPO emerges as the only method that is computationally efficient while offering consistent improvements over the Anthropic HH preferred completions and matches or surpasses the resource-intensive Best of 128 approach. Lastly, Figure~\ref{fig:dialogue-main} highlights that DPO rapidly attains its peak effectiveness during training.

\subsection{Generalization to a new input distribution}

To deepen our understanding of how PPO and DPO perform under distributional changes, we examine their generalization capabilities by applying the trained policies from our Reddit TL;DR summarization study to a new domain: news articles from the test set of the CNN/DailyMail dataset \citep{nallapati-etal-2016-abstractive}. The same optimal sampling temperatures identified for TL;DR (0 and 0.25) are utilized in this assessment. The corresponding outcomes are summarized in Table~\ref{tab:ood}. To measure summary quality, we calculated the win rate of GPT-4 against the ground-truth dataset summaries, using a GPT-4 (C) prompt analogous to that in the Reddit TL;DR setup, but substituting the phrase ``forum post'' with ``news article''. Within this distinct distribution, DPO consistently achieves substantially better performance than PPO. These findings offer preliminary evidence that DPO’s ability to generalize is comparable to PPO, despite DPO not leveraging the extra unlabeled Reddit TL;DR prompts that are accessible to PPO.

\subsection{Validating GPT-4 judgments with human judgments}
\label{sec:human-judgments}

To assess the validity of GPT-4 as an automatic evaluator, we ran a human assessment based on summaries from the TL;DR experiment, utilizing two separate prompts for GPT-4. The first, \textbf{GPT-4 (S)} (simple), asks directly which summary more effectively captures the post’s key points. The second, \textbf{GPT-4 (C)} (concise), also queries conciseness, included after observing that GPT-4 using the \textbf{GPT-4 (S)} prompt tends to favor longer, potentially redundant summaries more than human raters do. The precise prompt wording can be found in Appendix~\ref{app:prompts}. For the evaluation, we selected three settings to cover a range of sample qualities: the top-performing approach (DPO at temp. 0.25), the lowest (PPO at temp. 1.0), and an intermediate baseline (SFT at temp. 0.25); each is compared to PPO with greedy sampling, which performed best at its optimal temperature. Analysis reveals that, regardless of the prompt, GPT-4’s preferences align with human rater choices about as frequently as the raters agree with each other—supporting its use as a proxy for human evaluation (due to the limited number of raters, multiple human judgments were collected only for the DPO and PPO-1 matchups). Overall, \textbf{GPT-4 (C)} better mirrors human decisions in terms of win rates, motivating our use of this prompt for the primary experiments in Section~\ref{sec:dpo-real-datasets}. Details of the human evaluation methodology, such as the user interface and volunteer participant information, are provided in Appendix~\ref{app:human-study}.

\section{Discussion}
Preference-based learning presents a robust and extensible approach for aligning and enhancing the abilities of language models. In this work, we have proposed DPO, a streamlined methodology that enables training language models using preference data, circumventing the need for reinforcement learning. DPO eschews the conventional approach of reframing the preference learning task as a reinforcement learning (RL) problem merely to accommodate existing RL algorithms. Instead, it establishes a principled connection between language model policies and reward functions, facilitating the direct optimization of models to reflect human preferences via a standard cross-entropy objective—thus eliminating the need for reinforcement learning while retaining generality. With minimal hyperparameter adjustment, DPO attains performance on par with or superior to leading RLHF techniques, such as those leveraging PPO; consequently, DPO significantly lowers the practical obstacles to developing language models guided by human feedback.

\textbf{Limitations \& Future Work.} Our findings also point toward several compelling avenues for subsequent research. One open question concerns the ability of policies trained with DPO to generalize to distributions beyond the training data, relative to models optimized with explicit reward signals. Preliminary experiments indicate that DPO-trained policies exhibit generalization behavior similar to PPO-based approaches, though a more systematic exploration is warranted. For instance, it remains to be explored whether employing DPO-driven self-labeling can also exploit unannotated prompts effectively. Another aspect for investigation is the manifestation of reward over-optimization within the DPO framework and whether the minor performance reduction seen in Figure~\ref{fig:dialogue-main}-right exemplifies such effects. Additionally, while our current study evaluates models up to 6B parameters, scaling DPO to state-of-the-art, much larger models represents an intriguing future direction. With respect to assessment, we observe that GPT-4-based win rate evaluations are prompt-dependent; further research could address how best to elicit reliable judgments from automated evaluators. Lastly, there is considerable potential for DPO to be extended to the training of generative models beyond language, across diverse modalities.